\begin{claim}{B.29}{Proved}{Complete executed punctured genus-two catalogues}
The extended mapping class groups of types $(2,1)$ and $(2,2)$ have
respectively $32$ and $121$ conjugacy classes of finite subgroups.
Their positive subgroups form $10$ and $23$ of these classes.
The files \texttt{type\_2\_\nu\_extended\_catalogue.csv} and
\texttt{type\_2\_\nu\_extended\_sections.json}, for $\nu=1,2$,
list every class with its exact group, puncture orbits and stabilizers,
all section classes and their incidences, $\Phi$, its group action
and stabilizers, centralizers, and all section and puncture multipliers.
\end{claim}
\begin{proof}
Fill the punctures of a finite action on the genus-two surface.
The resulting finite subgroup in the closed mapping class group is
faithful: a finite homeomorphism isotopic to the identity on a closed
hyperbolic surface is the identity, by B.6.  Thus B.28 supplies the
complete list of closed actions that must be decorated.
For such an action $K$, keep the notation
$E=K\langle\iota\rangle$, $H=E/\langle\iota\rangle$, and
$Q=S^2/H$ from that proof.

The strata relevant to punctures are intrinsic to $K$.
Let $\mathcal V_K$ be the set of points whose stabilizer in $K$
contains a positive nonidentity element.  Take its $K$-orbits.
Cut the fixed circles of the reflections in $K$ at
$\mathcal V_K$ and take the $K$-orbits of the resulting open
intervals or uncut circles.  These are the mirror strata.
There is one connected regular interior in the quotient $S_2/K$.
A vertex orbit is selected zero or one time; a mirror stratum
receives any nonnegative number of puncture orbits; so does the
regular interior.  Their weights are, respectively,
\[
       |Kp|,\qquad |K|/2,\qquad |K|.                  \tag{B.29.1}
\]
The weighted sum must equal $\nu$.
Points fixed by $\iota$ but having no positive stabilizer in $K$
are ordinary points of these $K$-strata.  They introduce no extra
puncture label.  In particular punctures need not be preserved
by $\iota$ when $K$ is a complement in $E$.

We prove the normalizer permutations used to identify these
decorations.  Every mapping class normalizing $K$ normalizes $E$,
since $\iota$ is central in $\mathcal T_{2,0}$.  At the homeomorphism
level it may be represented by a normalizer of the chosen finite
action $E$ without changing its action on the intrinsic $K$-strata.
Indeed compare the two commuting involutions $\iota$ and
$h\iota h^{-1}$ after conjugating $K$ by $h$.  The finite actions of
the abstract extension $K\langle\iota\rangle$ have the same mapping
classes.  Uniqueness of realization in C.6 gives a conjugating
homeomorphism isotopic to the identity and commuting with $K$.
C.7 gives an isotopy through such commuting homeomorphisms.
It fixes the isolated $K$-fixed points individually and preserves
every reflection component.  Composing with it makes $h$ normalize
the chosen $E$ and leaves the required stratum permutation unchanged.
Thus it suffices, and is necessary, to compute the action of the
normalizer of $E$ that preserves $K$.

The full-group case $K=E$ is direct.  A cone in $Q$ has the color
consisting of its original order and its branch flag.
An ordinary branch orbit has its own color $(1,1)$.
The normalizer permits all permutations of identical interior
cone colors and exactly the dihedral symmetries of the expanded
colored boundary list of B.28.  Each is realized by a quotient
orbifold homeomorphism and hence lifts.  Pure interior braids
conjugate local generators, so they leave the corresponding
$E$-orbits unchanged.  These observations prove that this finite
permutation group is the whole normalizer action on the strata.

Suppose now $K=\ker\chi$ with $\chi(z)=1$.
For a local stabilizer $L\subset E$, its point orbit is $E/L$.
It is a single $K$-orbit when $\chi(L)=\mathbf F_2$.
When $\chi(L)=0$, it consists of two $K$-orbits, labelled
\[
                       b=\chi(g),\qquad gL\in E/L.    \tag{B.29.2}
\]
This is well defined precisely in the latter case.
It gives the exact orbit comparison, including the ambiguity in
choosing a lift of a quotient homeomorphism: composing with $z$
simultaneously changes every available sheet label $b$ to $b+1$.
The code includes this deck permutation.

An interior cone with generator $x$ contributes an intrinsic
$K$-vertex only when $K\cap\langle x\rangle$ is nontrivial.
For cone permutations, retain exactly the equal-color permutations
that also preserve the character values on their local generators.
On a sphere quotient, every split cone pair has an additional
independent interchange.  Here is the explicit lift proving that
assertion.  Choose a branch local generator $y$ with $\chi(y)=1$;
one exists because the branch decoration is nonempty.
An ordinary branch generator is $z$.  A branched old cone has
$y^m=z$ with $m$ odd whenever a character with $\chi(z)=1$ exists,
and therefore also has $\chi(y)=1$.
For a split cone generator $x$, we have $\chi(x)=0$.
A pure braid about these two quotient points acts by
\[
       x\longmapsto(xy)x(xy)^{-1},
       \qquad y\longmapsto xyx^{-1}.                 \tag{B.29.3}
\]
The other local generators are unchanged.
The character is preserved.  The conjugator for $x$ has character
one, so (B.29.2) interchanges its two $K$-orbits; the conjugator for
$y$ has character zero.  Thus it leaves the other split pairs
unchanged.  This proves the claimed independent interchanges.
Together with cone permutations they generate the whole stratum
action: pure braids generate the remaining mapping classes of a
punctured sphere, and their action on the two-element fibers can
only be these interchanges.

For a disc quotient there is at most one old interior cone.
An ordinary interior branch point permits exactly the same
interchange (B.29.3) of that split cone pair, supported away from
the boundary.  When there is no such interior branch point, the
disc has at most one interior marked cone; its pure braid group
introduces no additional permutation.  Boundary permutations are
computed as follows.  Let $L$ be the length of the expanded corner
list, write $R_{j+L}=eR_je^{-1}$, and put
$\epsilon=\chi(e)$.  A colored cyclic shift sends a corner index
$j$ to $s+j$ and an edge index $i$ to $s+i$.
A reversal sends them to $s-j-1$ and $s-i$, respectively.
For either image write the resulting integer as $qL+t$,
$0\le t<L$.  Its local generator is conjugated by $e^q$;
hence the sheet label changes by
\[
                          b\longmapsto b+q\epsilon.   \tag{B.29.4}
\]
Retain precisely the colored symmetries that preserve the reflection
character values, and include the global deck interchange.
These formulas give both vertex and mirror-component permutations.
For a cornerless boundary the same rule is its circular covering
with monodromy $\chi(e)$; the lifted mirror graph records whether
the two local sheet labels belong to one or two components.

There are no omitted boundary permutations.  An orbifold
homeomorphism of the disc preserves its cyclic boundary order
up to reversal.  After its colored dihedral permutation has been
removed, its remaining mapping class is an interior braid or a
boundary twist.  An interior braid has the local conjugators
already described and fixes the boundary.  A boundary twist has
the conjugator $e$, already represented by (B.29.4) and the deck
permutation.  Inner automorphisms coming from $K$ leave $K$-orbits
unchanged.  The projective-plane cases in B.28 have no complements;
their full-group decorations require only permutations of identical
interior cone colors.  Thus every case has an explicit complete
finite normalizer permutation group.

We spell out how this calculation handles smooth points of the
$K$-mirror locus that were corners for $E$.  The code starts with
the exact finite mirror coset graph of the full action $E$.
It finds the vertices with positive stabilizer in $K$, removes
only these vertices to form the open mirror strata, and computes
the fixed circles of each actual reflection in $K$.
Thus an $E$-corner that has only a reflection stabilizer in $K$
lies inside a single $K$-mirror stratum.  The sheet and boundary
formulas above are applied to an ordinary edge sample in each
such component.  They determine its whole component by
connectedness.  This is also a direct check that a puncture can
move through a noncritical Weierstrass point whenever the
$K$-stratum allows it.

Enumerate all vectors with the weights (B.29.1) summing to
$\nu=1$ or $2$, and divide by the proved finite permutation group.
These are precisely the topological action classes with punctures.
Indeed unlabelled finite subsets of an interval or circle with
the same cardinality are related by a homeomorphism of prescribed
degree, and the regular interior configuration space is connected.
Extend such moves by small discs and strips in the quotient and
lift the isotopies through the finite orbifold cover.  Conversely
an equivariant homeomorphism preserves the intrinsic strata,
their weights, and a stratum permutation in the normalizer just
computed.  C.5--C.6 then identify these action classes with the
finite-subgroup conjugacy classes in the punctured mapping class
group.

The executed program \texttt{catalogue\_genus2\_punctures.py}
returns $32$ and $121$ classes.  It places the punctures on the
exact coset graph, with one midpoint inserted in every edge before
any vertex is removed.  This subdivision guarantees that even
two punctures at formerly adjacent samples leave the intervening
open interval represented.  The program checks that an element
stabilizing an edge fixes both its endpoints, so the inserted
midpoint has exactly the edge stabilizer and creates no artificial
critical point.  It records regular puncture orbits by the free
$K$-set $K$, with trivial stabilizer.

All positive section classes from B.28 are retained exactly when
their fixed point has not been removed and their image subgroup
lies in $K$.  For every reflection, remove the punctures from its
actual fixed-circle graph and recompute the resulting ovals and
lines.  These operations give all partial-section classes by
B.3--B.4.  Their inclusion poset, maximal classes, $K$-action,
stabilizers and centralizers are then recomputed.
Local multipliers at punctures are retained from the positive
local generators or assigned $1/2$ to a reflection.
The positive portions contain $10$ and $23$ classes and agree
with the independent primary-source calculation in B.33 on the
full multiset of group order, full sections, partial sections,
$|\Phi|$, and $|\Phi/K|$.  The negative involutions provide
another independent check: there are $4$ classes with one
puncture and $11$ with two, exactly the decorations of the five
closed real types in B.16.
\end{proof}
